% =====================================================================
% Análise heurística, problemas de usabilidade e correções
% Colocar no chap3.tex, antes de \section{Vídeo de Demonstração}:
%     \input{heuristica}
% Pacotes necessários (já tens): booktabs, longtable, array, float, graphicx
% =====================================================================
\chapter{Análise Heurística}
% --- Quadro de um problema de usabilidade (6 campos) ---
\newcommand{\problema}[6]{%
  {\small
  \begin{longtable}{>{\bfseries\raggedright\arraybackslash}p{3.3cm}
                    >{\raggedright\arraybackslash}p{\dimexpr\textwidth-3.3cm-4\tabcolsep\relax}}
    \toprule
    Ecrã ou percurso & #1 \\
    \midrule
    Descrição & #2 \\
    \midrule
    Heurística & #3 \\
    \midrule
    Consequências & #4 \\
    \midrule
    Gravidade & #5 \\
    \midrule
    Proposta de correção & #6 \\
    \bottomrule
  \end{longtable}}
}

% --- Figura de comparação (se a imagem não existir, mostra uma caixa) ---
% Uso: \figcomparacao{ficheiro.png}{legenda}{label}
\newcommand{\figcomparacao}[3]{%
  \begin{figure}[H]
    \centering
    \IfFileExists{images/#1}%
      {\includegraphics[width=\textwidth]{images/#1}}%
      {\fbox{\parbox{0.9\textwidth}{\centering\vspace{2cm}\textbf{Inserir imagem: images/#1}\vspace{2cm}}}}
    \caption{#2}
    \label{#3}
  \end{figure}
}

% =====================================================================
\section{Análise Heurística Simplificada}
\label{sec:heuristica}

A versão anterior do protótipo foi avaliada percorrendo as tarefas principais dos três perfis,
com especial atenção ao percurso do Cliente (``Garantir mesa para jantar''), e comparando cada
ecrã com as heurísticas de usabilidade abordadas nas aulas. São apresentadas apenas as heurísticas relevantes para o protótipo, com os aspetos positivos e os problemas
encontrados.

{\small
\begin{longtable}{>{\raggedright\arraybackslash}p{3cm}
                  >{\raggedright\arraybackslash}p{5.6cm}
                  >{\raggedright\arraybackslash}p{5.6cm}}
  \caption{Análise heurística da versão anterior do protótipo}\label{tab:heuristicas}\\
  \toprule
  \textbf{Heurística} & \textbf{Aspetos positivos} & \textbf{Problemas} \\
  \midrule
  \endfirsthead
  \toprule
  \textbf{Heurística} & \textbf{Aspetos positivos} & \textbf{Problemas} \\
  \midrule
  \endhead
  \bottomrule
  \endfoot

  Visibilidade do estado do sistema
  & O estado das reservas e da fila era sempre visível, através de etiquetas com texto e ícone
    (``Confirmada'', ``Cancelada'', ``Não compareceu''), e o ecrã ``A confirmar\ldots''
    informava que o pedido estava em curso.
  & Os filtros do ecrã de início e os seletores de hora e de número de pessoas não reagiam ao
    toque, pelo que o utilizador não sabia se a escolha tinha ficado registada (problema~2).
    Em vários ecrãs, o conteúdo ficava cortado no fundo, sem indicação de que havia mais
    informação (problema~1). \\
  \addlinespace

  Correspondência entre o sistema e o mundo real
  & A linguagem usa termos do dia a dia da restauração (``Reservar mesa'', ``Entrar na fila de
    espera'', ``Mesa quase pronta''). As datas, horas e números de pessoas são apresentados
    como o utilizador os diz.
  & \emph{Problema menor:} os cartões de restaurante mostravam um ícone genérico em vez de uma
    fotografia, o que dificultava reconhecer o restaurante. \\
  \addlinespace

  Controlo e liberdade do utilizador
  & As ações destrutivas (cancelar uma reserva, sair da fila) já tinham um diálogo com a opção
    de manter.
  & A barra inferior não ficava fixa e o botão voltar tinha pouco contraste (problema~3).
    Faltavam também alguns ecrãs, como a confirmação de envio da ligação de recuperação da
    palavra-passe e a eliminação da conta, o que deixava percursos sem saída. \\
  \addlinespace

  Consistência e utilização de convenções
  & A aplicação segue convenções habituais em aplicações móveis: barra de navegação inferior,
    botão voltar no canto superior esquerdo, um botão principal por ecrã.
  & Os elementos não estavam construídos como componentes reutilizáveis, o que originava
    pequenas diferenças de alinhamento e de centragem de ícones e texto entre ecrãs. \\
  \addlinespace

  Prevenção de erros
  & As horas esgotadas apareciam assinaladas antes de o utilizador as escolher.
  & Muitos elementos interativos tinham áreas de toque abaixo do mínimo recomendado, o que
    favorecia toques falhados ou no elemento errado (problema~4). \\
  \addlinespace

  Reconhecimento em vez de memorização
  & Os separadores da barra inferior combinam ícone e texto.
  & Os campos de texto tinham a mesma cor dos cartões informativos, pelo que não se
    reconheciam de imediato como zonas de escrita (problema~5). \\
  \addlinespace

  Flexibilidade e eficiência de utilização
  & O percurso principal é curto (três toques desde o ecrã de início até à reserva).
  & Os atalhos de pesquisa (filtros) não funcionavam, o que obrigava a percorrer toda a lista. \\
  \addlinespace

  Design estético e minimalista
  & Os ecrãs mostram apenas a informação necessária.
  & \emph{Problema estético:} o espaçamento vertical entre grupos de conteúdo era reduzido, o
    que tornava alguns ecrãs densos. \\
  \addlinespace

  Reconhecimento e recuperação de erros
  & Os erros de formulário (por exemplo, palavra-passe incorreta) eram apresentados junto ao
    campo, com ícone e uma mensagem que indicava como corrigir.
  & Depois de pedir a recuperação da palavra-passe, o utilizador não recebia confirmação. \\
  \addlinespace

  Ajuda e informação de apoio
  & A aplicação é simples e quase não precisa de ajuda; os campos mais exigentes têm texto de
    apoio (``Mínimo de 8 caracteres'').
  & \emph{Problema estético:} a ligação ``Política de privacidade'' no registo não levava a
    nenhum ecrã. \\
\end{longtable}}


% =====================================================================
\section{Problemas de Usabilidade}
\label{sec:problemas}

Cada problema é descrito com o ecrã ou percurso afetado, a heurística violada, as consequências
para o utilizador, a gravidade (escala de 0 a 4 de Nielsen~\cite{nielsen}) e a correção proposta.

\subsection{Problema 1: conteúdo cortado nos ecrãs longos}

\problema
  {Painel do serviço (Gestor) e, em geral, todos os ecrãs com muito conteúdo, como o Início e
   As minhas reservas.}
  {Os ecrãs não tinham \textit{scroll}, por isso todo o conteúdo tinha de caber na altura do
   telemóvel. Para isso, os elementos ficavam comprimidos: no Painel do serviço só cabiam três
   indicadores, apertados, e o texto do botão ``Não compareceu'' quase tocava na borda. As listas
   também não podiam crescer além do que cabia no ecrã.}
  {Visibilidade do estado do sistema.}
  {Informação importante ficava de fora ou difícil de ler, e as listas não podiam mostrar mais
   restaurantes ou reservas sem cortar conteúdo.}
  {3 \textperiodcentered{} Problema grave.}
  {Ativar o \textit{scroll} vertical em todos os ecrãs cujo conteúdo excede a altura do ecrã,
   mantendo fixos os elementos de navegação.}

\subsection{Problema 2: filtros e seletores sem resposta}

\problema
  {Início (filtros) e Reservar mesa (seleção de dia, hora e número de pessoas), no percurso
   principal do Cliente.}
  {Os filtros do ecrã de início (``Hoje, 20:00'', ``2 pessoas'', ``Cozinha'') e os seletores do
   ecrã de reserva não reagiam ao toque. A opção selecionada não mudava e a lista de restaurantes
   mantinha-se igual.}
  {Visibilidade do estado do sistema; Flexibilidade e eficiência de utilização.}
  {O utilizador não sabia se a escolha tinha sido registada, podia julgar que a aplicação estava
   avariada e não conseguia reduzir a lista ao tipo de cozinha pretendido.}
  {3 \textperiodcentered{} Problema grave.}
  {Tornar os seletores interativos, com estados visíveis (predefinido, selecionado, esgotado), e
   substituir os filtros por filtros de cozinha que alteram a lista apresentada.}

\subsection{Problema 3: navegação pouco persistente e botão voltar pouco visível}

\problema
  {Todos os ecrãs com barra inferior (Início, As minhas reservas, A sua fila) e todos os ecrãs
   com cabeçalho (por exemplo, Reservar mesa).}
  {A barra de navegação inferior não ficava fixa no fundo do ecrã. O botão voltar era um círculo
   bege sobre fundo creme, com pouco contraste.}
  {Controlo e liberdade do utilizador.}
  {O utilizador podia perder o acesso à navegação principal e não identificar de imediato como
   regressar ao ecrã anterior, sentindo-se ``preso''.}
  {3 \textperiodcentered{} Problema grave.}
  {Fixar a barra inferior e a barra de estado em todos os ecrãs e redesenhar o botão voltar com
   maior contraste.}

\subsection{Problema 4: áreas de toque pequenas}

\problema
  {Vários ecrãs, sobretudo Reservar mesa (opções de hora), Início (filtros) e os cabeçalhos
   (botão voltar).}
  {As ligações eram zonas desenhadas à parte, por cima de imagens importadas, e várias eram mais
   pequenas do que o botão visível, abaixo dos 44 $\times$ 44 px recomendados para ecrãs
   táteis.}
  {Prevenção de erros.}
  {Toques falhados ou no elemento ao lado, por exemplo escolher a hora errada, o que obriga o
   utilizador a repetir ações. É especialmente penalizador para quem usa o telemóvel a caminho do
   restaurante, numa só mão.}
  {3 \textperiodcentered{} Problema grave.}
  {Definir uma área de toque mínima de 44 $\times$ 44 px para todos os elementos interativos e
   aumentar o espaçamento entre eles.}

\subsection{Problema 5: campos de texto pouco reconhecíveis}

\problema
  {Iniciar sessão, Criar conta e Pesquisa, no início do percurso principal.}
  {Os campos de texto tinham a mesma cor de fundo (bege) dos cartões informativos e só se
   distinguiam deles pela borda.}
  {Reconhecimento em vez de memorização; Consistência e utilização de convenções.}
  {O utilizador demorava mais a perceber onde podia escrever, sobretudo nos formulários com
   vários campos.}
  {2 \textperiodcentered{} Problema menor.}
  {Dar aos campos de texto uma cor própria, mais clara que a dos cartões, mantendo a borda, e
   aplicá-la de forma igual em todos os campos da aplicação.}

% =====================================================================
\section{Correções e Melhorias}
\label{sec:correcoes}

Todas as correções foram feitas nos componentes do \textit{design system} do Figma. Assim, cada
alteração foi aplicada uma só vez e propagou-se automaticamente a todos os ecrãs que usam esse
componente.

\subsection{Correção 1: \textit{scroll} nos ecrãs longos}

\begin{itemize}
  \item \textbf{Problema resolvido:} problema~1.
  \item \textbf{Versão anterior:} ecrã ``Painel do serviço'' (Gestor), com três indicadores
        apertados e o botão ``Não compareceu'' no limite.
  \item \textbf{Versão corrigida:} ecrã ``31 \textperiodcentered{} Painel do serviço'', com
        quatro indicadores e \textit{scroll} vertical.
  \item \textbf{Justificação:} todos os ecrãs cujo conteúdo excede a altura do telemóvel
        passaram a ter \textit{scroll} vertical. A barra de estado e a barra inferior ficam
        fixas, para que o utilizador veja todo o conteúdo sem perder a navegação. O espaço
        inferior do conteúdo foi aumentado, para que o último elemento nunca fique escondido
        atrás da barra. Deixando de estar limitado à altura do ecrã, o painel passou a mostrar
        quatro indicadores, com espaço entre eles.
\end{itemize}

\figcomparacao{comparacao1.png}{Comparação 1: Painel do serviço antes e depois da correção}{fig:comp1}

\subsection{Correção 2: filtros e seletores interativos}

\begin{itemize}
  \item \textbf{Problema resolvido:} problema~2.
  \item \textbf{Versão anterior:} ecrãs ``Início'' e ``Reservar mesa'', com filtros e seletores
        estáticos.
  \item \textbf{Versão corrigida:} ecrãs ``06 \textperiodcentered{} Início'' e
        ``06c \textperiodcentered{} Início (Japonesa)'', com a lista filtrada, e
        ``09 \textperiodcentered{} Reservar mesa'', com os seletores interativos.
  \item \textbf{Justificação:}
    \begin{itemize}
      \item \textbf{Filtros:} passaram a ser por tipo de cozinha (Todas, Portuguesa, Japonesa,
            Petiscos); ao tocar num filtro, a lista mostra só os restaurantes dessa cozinha e o
            filtro fica assinalado;
      \item \textbf{Seletores de dia, hora e número de pessoas:} mudam de estado ao toque, com
            indicação clara da opção selecionada;
      \item \textbf{Hora esgotada:} aparece tracejada e com a palavra ``esgotado''; ao tocar
            nela, o utilizador vê uma explicação com horas alternativas.
    \end{itemize}
\end{itemize}

\figcomparacao{comparacao2.png}{Comparação 2: filtros e seletores antes e depois da correção}{fig:comp2}

\subsection{Correção 3: navegação fixa e botão voltar mais visível}

\begin{itemize}
  \item \textbf{Problema resolvido:} problema~3.
  \item \textbf{Versão anterior:} cabeçalho de ``Reservar mesa'', com o botão voltar bege, e
        ``Início'', com a barra inferior não fixa.
  \item \textbf{Versão corrigida:} os mesmos ecrãs, com o botão voltar verde e a barra inferior
        fixa.
  \item \textbf{Justificação:}
    \begin{itemize}
      \item \textbf{Barra inferior:} fica sempre visível, com o separador ativo assinalado por um
            indicador e texto em negrito;
      \item \textbf{Botão voltar:} passou a ser um círculo verde com seta creme, com um contraste
            de 5,8:1, no mesmo sítio em todos os ecrãs;
      \item \textbf{Ecrãs em falta:} acrescentaram-se os ecrãs que faltavam (confirmação do envio
            da ligação de recuperação, política de privacidade, alteração da fotografia e
            eliminação da conta), para que nenhum percurso fique sem saída.
    \end{itemize}
\end{itemize}

\figcomparacao{comparacao3.png}{Comparação 3: navegação e botão voltar antes e depois da correção}{fig:comp3}

\subsection{Correção 4: áreas de toque com o mínimo recomendado}

\begin{itemize}
  \item \textbf{Problema resolvido:} problema~4.
  \item \textbf{Versão anterior:} ecrã ``Reservar mesa'', com as ligações em zonas desenhadas à
        parte, mais pequenas do que os botões.
  \item \textbf{Versão corrigida:} componentes com as novas dimensões, por exemplo no ecrã
        ``09 \textperiodcentered{} Reservar mesa''.
  \item \textbf{Justificação:} todos os elementos interativos passaram a ter pelo menos
        44 $\times$ 44 px:
    \begin{itemize}
      \item botões: 52 px de altura;
      \item opções de hora e filtros: 44 px;
      \item botão voltar, avatar e sino: 48 $\times$ 48 px;
      \item itens da barra inferior: 97,5 $\times$ 72 px.
    \end{itemize}
    O espaço entre grupos de conteúdo passou para 24 px, o que também reduz toques no elemento
    errado. As dimensões foram verificadas em todos os elementos com ligação no protótipo.
\end{itemize}

\figcomparacao{comparacao4.png}{Comparação 4: áreas de toque antes e depois da correção}{fig:comp4}

\subsection{Correção 5: campos de texto com cor própria}

\begin{itemize}
  \item \textbf{Problema resolvido:} problema~5.
  \item \textbf{Versão anterior:} ecrã ``Iniciar sessão'', com os campos em bege, iguais aos
        cartões.
  \item \textbf{Versão corrigida:} ecrã ``02 \textperiodcentered{} Iniciar sessão'', com os
        campos em bege claro (\autoref{fig:comp5}).
  \item \textbf{Justificação:} acrescentou-se à paleta uma nova cor, o bege claro
        (\texttt{\#F3EAD7}), usada em todos os campos de texto e na pesquisa, mantendo a borda
        em areia. Os campos distinguem-se assim dos cartões informativos. O texto escrito mantém
        um contraste de 9,8:1, acima do nível AA da WCAG~\cite{wcag}.
\end{itemize}

\figcomparacao{comparacao5.png}{Comparação 5: campos de texto antes e depois da correção}{fig:comp5}

\subsection{Outras melhorias}

Além dos cinco problemas descritos, a análise motivou outras alterações:

\begin{itemize}
  \item \textbf{Fotografias:} os cartões e os detalhes de restaurante passaram a mostrar
        fotografias;
  \item \textbf{Ecrã de abertura:} passou a existir em todos os percursos, incluindo os do Gestor
        e do Administrador;
  \item \textbf{Organização no Figma:} os componentes ficaram agrupados num \textit{design
        system}, com nomes claros e variantes de estado, e os ecrãs em secções por percurso;
  \item \textbf{Fidelidade ao sistema visual:} confirmou-se que as cores, a tipografia, a grelha
        e as margens do protótipo correspondem ao definido na \autoref{sec:sistema-visual}.
\end{itemize}